\hypertarget{appendix-8404d0b1a3e7}{}\label{referencestart:appendix-8404d0b1a3e7}
\section{Sparse-set storage}\label{appendix-8404d0b1a3e7-sparse-set-storage}

\hypertarget{appendix-8404d0b1a3e7-sparse-set-storage}{}

A sparse set answers a common ECS question: how do we keep component values densely packed while entities live in a much larger, sparse identifier space?

Use two index structures:

\begin{itemize}
\tightlist
\item
  dense: a compact sequence of live entities, often parallel to a compact component-value array;
\item
  sparse: indexed by entity index and storing where that entity should appear in dense.
\end{itemize}

If sparse{[}e.index{]} says dense slot d, membership is not established until dense{[}d{]} is checked against the complete entity identity. This matters when entity indices are reused with newer generations. An old handle and a new handle may share the same numeric index.

A common representation encodes absence as zero and stores dense\_index + 1 in sparse. Then:

contains(e) = 1. e.index is within sparse bounds; 2. sparse{[}e.index{]} is nonzero; 3. d = sparse{[}e.index{]} - 1 is within dense bounds; 4. dense\_entities{[}d{]} equals the full entity handle e.

Step 4 prevents a stale generation from inheriting the new occupant's component by index coincidence.

\subsection{Swap-and-pop removal}\label{appendix-8404d0b1a3e7-swap-and-pop-removal}

\hypertarget{appendix-8404d0b1a3e7-swap-and-pop-removal}{}

To erase dense slot d in O(1), move the last dense element into d and pop the last row. But the moved entity's sparse reverse mapping must be repaired.

If b is removed from {[}a,b,c{]}, c moves into b's old dense slot. The operation is therefore not just:

dense{[}d{]} = dense.back(); dense.pop\_back();

It also requires:

sparse{[}c.index{]} = d.

Forgetting that repair produces one of the nastiest ECS bugs: a membership query succeeds but returns the wrong component row later.

\subsection{Address stability}\label{appendix-8404d0b1a3e7-address-stability}

\hypertarget{appendix-8404d0b1a3e7-address-stability}{}

Swap-and-pop can invalidate a pointer even when the underlying vector does not allocate. The bytes at a saved address may now belong to another entity. Vector reallocation creates an additional invalidation mode. Therefore sparse-set component addresses should be treated as borrows with a clear mutation boundary.

Stable-address designs exist---paged pools, indirection, tombstones, node storage---but they trade locality, memory, branch cost, or iteration simplicity for stability.

\subsection{Complexity is not the whole workload}\label{appendix-8404d0b1a3e7-complexity-is-not-the-whole-workload}

\hypertarget{appendix-8404d0b1a3e7-complexity-is-not-the-whole-workload}{}

Sparse sets give expected O(1) membership, insertion and removal and dense O(n) iteration. But real performance depends on:

\begin{itemize}
\tightlist
\item
  size of the sparse index array;
\item
  cache behavior of random entity-index probes;
\item
  component payload size;
\item
  frequency of structural mutations;
\item
  query combinations across several component pools;
\item
  whether iteration follows the smallest pool or performs archetype-style grouping.
\end{itemize}

So ``O(1)'' does not settle the architecture decision.

\subsection{Executable reconstruction}\label{appendix-8404d0b1a3e7-executable-reconstruction}

\hypertarget{appendix-8404d0b1a3e7-executable-reconstruction}{}

tools/labs/ecs\_reconstruction\_lab.cpp implements a generational entity pool and a sparse-set component pool. Its tests remove a middle dense element and verify that the swap-moved entity still resolves correctly. It also checks stale-generation behavior.

Read the connected engine narrative at \hyperref[05-engine-i-ch014-let-identity-survive-changes-in-storage]{Continuous Notes/Game Engine - From First Frame to Living Worlds\#Let identity survive changes in storage}.

Related: \hyperref[appendix-fbfeb0430008]{Archetype ECS}, \href{../01-Bits-Programs-and-Reliable-Systems-P02/main.pdf\#nameddest=appendix-f1581457d322}{Object identity and generation handles (Bits, Programs, and Reliable Systems, Part 2, Reference notes: Object identity and generation handles)}, \href{../01-Bits-Programs-and-Reliable-Systems-P02/main.pdf\#nameddest=appendix-f9b8b9188c7f}{Borrow invalidation and logical identity (Bits, Programs, and Reliable Systems, Part 2, Reference notes: Borrow invalidation and logical identity)}.
